\begin{table}[tbp]
\centering
\small
\caption{118th Senate: networks simulated from the MPLE fits on the complete projection compared with the observed network.}
\label{tab:mle}
\begin{tabular}{llrrl}
\toprule
Model & statistic & observed & simulated mean & 95\% range \\
\midrule
M1 & edges & 4,290 & 4,283 & [4,165, 4,386] \\
M1 & mutual & 1,316 & 1,312 & [1,242, 1,370] \\
M1 & transitive triples & 103,456 & 67,856 & [62,701, 72,935] \\
M2 & edges & 4,290 & 4,436 & [4,315, 4,566] \\
M2 & mutual & 1,316 & 1,344 & [1,274, 1,426] \\
M2 & transitive triples & 103,456 & 79,598 & [73,157, 86,612] \\
M4c & edges & 4,290 & 3,163 & [3,026, 3,311] \\
M4c & mutual & 1,316 & 720 & [656, 780] \\
M4c & transitive triples & 103,456 & 29,846 & [26,085, 34,275] \\
M5 & edges & 4,290 & 3,751 & [3,430, 4,030] \\
M5 & mutual & 1,316 & 978 & [843, 1,102] \\
M5 & transitive triples & 103,456 & 53,991 & [41,831, 67,353] \\
S0 & edges & 4,290 & 4,297 & [4,199, 4,378] \\
S0 & mutual & 1,316 & 1,317 & [1,254, 1,365] \\
S0 & transitive triples & 103,456 & 100,228 & [93,028, 106,242] \\
S0 & cyclic triples & 31,730 & 30,613 & [28,240, 32,517] \\
S1 & edges & 4,290 & 4,276 & [4,183, 4,364] \\
S1 & mutual & 1,316 & 1,311 & [1,261, 1,366] \\
S1 & transitive triples & 103,456 & 100,869 & [94,869, 106,942] \\
S1 & cyclic triples & 31,730 & 30,842 & [28,853, 32,729] \\
\bottomrule
\end{tabular}
\par\vspace{2pt}\footnotesize One hundred networks per model simulated with the statnet \texttt{ergm} package (burn-in 200,000, interval 10,000) from the MPLE coefficients, which agree with the Python MPLE to four decimal places. M models use geometrically weighted degree terms; S models use sender and receiver effects.
\end{table}